\documentclass[11pt,a4paper]{article}
\usepackage[utf8]{inputenc}
\usepackage[T1]{fontenc}
\IfFileExists{french.ldf}{\usepackage[french]{babel}}{\usepackage[english]{babel}}
\usepackage{amsmath,amssymb,amsthm}
\usepackage{geometry}\geometry{margin=2.3cm}
\usepackage{listings}
\usepackage{xcolor}
\usepackage{enumitem}
\usepackage{fancyhdr}
\usepackage{lmodern}
\usepackage{titlesec}
\usepackage{hyperref}
\hypersetup{colorlinks=true, urlcolor=[RGB]{2,101,115}, linkcolor=black}
\definecolor{codebg}{RGB}{247,249,251}
\definecolor{kw}{RGB}{175,45,45}
\definecolor{cmt}{RGB}{110,120,130}
\definecolor{str}{RGB}{20,110,60}
\lstset{language=Python, backgroundcolor=\color{codebg}, basicstyle=\ttfamily\small,
  keywordstyle=\color{kw}\bfseries, commentstyle=\color{cmt}\itshape, stringstyle=\color{str},
  numbers=left, numberstyle=\tiny\color{cmt}, numbersep=8pt, frame=single, rulecolor=\color{gray!40},
  showstringspaces=false, breaklines=true, columns=fullflexible, extendedchars=true, inputencoding=utf8,
  literate={é}{{\'e}}1 {è}{{\`e}}1 {à}{{\`a}}1 {ê}{{\^e}}1 {î}{{\^i}}1 {ç}{{\c c}}1 {ô}{{\^o}}1 {û}{{\^u}}1 {ù}{{\`u}}1 {É}{{\'E}}1}
\newcounter{qcnt}
\newcommand{\Q}{\medskip\par\stepcounter{qcnt}\noindent\textbf{Question~\theqcnt.}~}
\newcommand{\code}[1]{\lstinline!#1!}
\pagestyle{fancy}\fancyhf{}
\lhead{\small Préparation au CNC 2026 --- Informatique MP}
\rhead{\small Énoncé}
\cfoot{\thepage}
\renewcommand{\headrulewidth}{0.4pt}
\titleformat*{\section}{\large\bfseries\scshape}
\titleformat*{\subsection}{\normalsize\bfseries}
\begin{document}

\begin{center}
{\large\bfseries Préparation au Concours National Commun 2026 --- Informatique MP}\\[8pt]
\rule{0.9\linewidth}{0.5pt}\\[6pt]
{\Large\bfseries Épreuve d'Informatique --- Énoncé}\\[4pt]
{\bfseries Filière MP --- Durée : 4 heures --- Quatre problèmes indépendants}\\[6pt]
{\itshape Autour de la blockchain Ethereum : algorithmique, arbres de Merkle, partition (programmation dynamique) et finance décentralisée}\\[6pt]
\rule{0.9\linewidth}{0.5pt}\\[8pt]
{\small Auteur~: \href{https://orcid.org/0000-0002-6108-7176}{\textbf{Pr Agrégé El Hadiq Zouhair}}}\\[3pt]
{\small \href{https://cpgeacademy.org}{\textbf{cpgeacademy.org}} \quad---\quad Tous droits réservés \textcopyright\ cpgeacademy.org}
\end{center}
\medskip
\begin{center}
\fcolorbox{gray}{gray!5}{\begin{minipage}{0.93\linewidth}\small
\textbf{Avis important.} Ce document est une \textbf{initiative personnelle} du professeur,
à but strictement pédagogique. Il ne s'agit \textbf{pas} d'un sujet officiel du Concours
National Commun~; il s'en inspire uniquement par la forme et l'esprit.
\smallskip
\textbf{Consignes.} L'épreuve comporte \textbf{quatre problèmes indépendants}. Les questions sont
numérotées \textbf{en continu} sur l'ensemble du sujet~; le corrigé reprend la même numérotation.
On pourra traiter les problèmes dans l'ordre de son choix.
\end{minipage}}
\end{center}
\bigskip

\setcounter{qcnt}{0}

\section*{Problème I --- Algorithmique de la blockchain Ethereum}
\noindent\emph{Présentation.} Ethereum est un ordinateur mondial décentralisé. Ce problème étudie, de manière \emph{largement
indépendante} partie par partie, quelques structures et algorithmes qui le sous-tendent~: la
mesure des ressources (le \emph{gaz}), l'encodage hexadécimal des adresses, les \emph{arbres de
Merkle}, la sélection aléatoire d'un validateur, et enfin un problème d'optimisation résolu par
\emph{programmation dynamique}.

\noindent\textbf{Conventions.}
\begin{itemize}[nosep]
  \item Le langage est Python 3. Les entiers Python sont de taille arbitraire.
  \item Pour une liste \code{L}, \code{len(L)} est sa longueur et \code{L[i]} son terme d'indice $i$ (à partir de $0$).
  \item $\lfloor x \rfloor$ désigne la partie entière de $x$. $\log_b$ est le logarithme en base $b$.
  \item On rappelle que \code{n // b} et \code{n \% b} sont le quotient et le reste de la division euclidienne.
\end{itemize}

\noindent Les questions sont numérotées \textbf{en continu} sur tout le sujet.

\medskip
\subsection*{Partie I --- Unités et compteur de gaz}

L'unité de compte d'Ethereum est l'\emph{ether}~; sa plus petite subdivision est le \emph{wei},
avec $1$ ether $= 10^{18}$ wei. Chaque instruction exécutée consomme une quantité fixe de
\emph{gaz}~; une transaction fixe un plafond de gaz au-delà duquel l'exécution est interrompue.

\Q Écrire une fonction \code{ether_vers_wei(e)} qui, pour un entier $e \ge 0$, renvoie le nombre
de wei correspondant à $e$ ethers.

\Q Écrire \code{gaz_restant(plafond, couts)} où \code{couts} est la liste des coûts (entiers
positifs) des instructions exécutées dans l'ordre. La fonction retranche les coûts un à un~; si,
avant une instruction, le gaz disponible est \emph{strictement} inférieur à son coût, l'exécution
s'arrête et la fonction renvoie la chaîne \code{"panne de gaz"}. Sinon elle renvoie le gaz restant.

\Q Un contrat exécute une boucle dont l'itération numéro $i$ (pour $i \ge 0$) coûte $g/2^{i}$
unités de gaz, où $g$ est un réel strictement positif. Montrer que le coût total
$\displaystyle C = \sum_{i=0}^{+\infty} \frac{g}{2^{i}}$ est fini et vaut exactement $2g$. En
déduire que, quel que soit le nombre d'itérations effectivement exécutées, le coût cumulé est
majoré par $2g$.

\subsection*{Partie II --- Encodage hexadécimal des adresses}

Une adresse Ethereum est un entier affiché en base $16$ (hexadécimal), avec l'alphabet
\code{"0123456789abcdef"}. On étudie l'encodage d'un entier en base $16$.

\Q Écrire \code{entier_vers_hex(n)} qui, pour un entier $n \ge 0$, renvoie la chaîne de ses
chiffres hexadécimaux (sans le préfixe \code{0x}), avec la convention \code{entier_vers_hex(0) == "0"}.

\Q Écrire \code{hex_vers_entier(s)} qui, pour une chaîne \code{s} de chiffres hexadécimaux
minuscules, renvoie l'entier qu'elle représente.

\Q Pour $n \ge 1$, montrer que le nombre de tours de la boucle de \code{entier_vers_hex(n)}
(donc le nombre de chiffres produits) est exactement $\lfloor \log_{16} n \rfloor + 1$. En déduire
la complexité temporelle de \code{entier_vers_hex} en fonction du nombre de chiffres.

\Q Justifier la \textbf{terminaison} de la boucle de \code{entier_vers_hex} à l'aide d'un
\emph{variant de boucle}, puis prouver la \textbf{correction}, c'est-à-dire que pour tout
$n \ge 0$, \code{hex_vers_entier(entier_vers_hex(n))} est égal à $n$.

\subsection*{Partie III --- Arbres de Merkle}

Un \emph{arbre de Merkle} résume un ensemble de données par un seul haché, la \emph{racine}. On
se donne une fonction de combinaison
$$\mathtt{parent(a, b)} = (31\,a + b) \bmod 1000003,$$
et les feuilles sont des entiers. La racine se calcule par niveaux~: à chaque niveau, on remplace
la liste courante par la liste des \code{parent} de couples consécutifs~; si la liste a une
longueur impaire, on \textbf{duplique son dernier élément} avant de l'appairer. On s'arrête
lorsqu'il ne reste qu'un élément.

\Q Écrire \code{racine_merkle(feuilles)} qui prend une liste non vide de feuilles et renvoie la
racine. Vérifier à la main que \code{racine_merkle([1, 2, 3, 4])} vaut $1120$.

\Q Soit $n$ le nombre de feuilles. Montrer que le nombre total d'appels à \code{parent} effectués
par \code{racine_merkle} est un $O(n)$. \emph{Indication}~: majorer la somme des tailles des
niveaux successifs.

\Q On veut prouver l'\emph{inclusion} d'une feuille sans révéler tout l'arbre. Écrire
\code{preuve_inclusion(feuilles, i)} qui renvoie la liste des couples $(\text{frère}, \text{côté})$
rencontrés en remontant de la feuille d'indice $i$ jusqu'à la racine, où le côté vaut
\code{'L'} ou \code{'R'} et indique si le frère est à gauche ou à droite. Écrire ensuite
\code{verifier(feuille, preuve, racine)} qui recalcule la racine à partir de la feuille et de la
preuve, et renvoie un booléen.

\Q Montrer par récurrence sur la hauteur de l'arbre que, si la fonction \code{parent} est
\emph{résistante aux collisions} (on ne sait pas produire $ (a,b) \ne (a',b')$ avec
$\mathtt{parent}(a,b) = \mathtt{parent}(a',b')$), alors une preuve d'inclusion acceptée par
\code{verifier} garantit que la feuille figure bien parmi les feuilles ayant produit la racine.
Quelle est la taille d'une preuve d'inclusion en fonction de $n$~?

\subsection*{Partie IV --- Sélection aléatoire d'un validateur}

À chaque créneau (\emph{slot}), un validateur donné est choisi comme proposeur avec probabilité
$p \in\, ]0,1[$, indépendamment des autres créneaux. On note $N$ le numéro (à partir de $1$) du
premier créneau où ce validateur est choisi.

\Q Donner, pour $k \ge 1$, la probabilité $\mathbb{P}(N = k)$. Reconnaître la loi de $N$ et en
déduire l'espérance $\mathbb{E}[N]$.

\Q Montrer que $\mathbb{P}(N > k) = (1-p)^{k}$, puis que $\mathbb{P}(N < +\infty) = 1$
(\emph{terminaison presque sûre}~: le validateur est choisi en temps fini avec probabilité $1$).

\Q Application numérique~: pour $p = 1/4$, donner $\mathbb{E}[N]$ et la probabilité
$\mathbb{P}(N \le 4)$ sous forme de fraction irréductible.

\subsection*{Partie V --- Partition équilibrée par programmation dynamique}

On dispose d'une liste $S$ d'entiers strictement positifs. Le \emph{problème de partition}
demande s'il existe un sous-ensemble de $S$ dont la somme vaut exactement la moitié de la somme
totale (les deux parts ont alors la même somme). C'est un problème classique lié aux preuves à
divulgation nulle.

\Q Montrer que si la somme totale $T = \sum S$ est impaire, la réponse est nécessairement
\emph{non}. On suppose désormais $T$ paire et on pose $C = T/2$.

\Q Le \emph{problème de décision} associé (« existe-t-il un sous-ensemble de somme $C$~? ») est
\textbf{NP}. Justifier en décrivant un \emph{certificat} et l'algorithme de vérification, ainsi
que sa complexité.

\Q On définit, pour $0 \le i \le \mathtt{len(S)}$ et $0 \le c \le C$, le booléen $A[i][c]$ = « il
existe un sous-ensemble des $i$ premiers éléments de $S$ de somme $c$ ». Donner la relation de
récurrence vérifiée par $A[i][c]$ et son initialisation.

\Q En déduire une fonction \code{peut_partitionner(S)} utilisant un tableau à une dimension
\code{dp} de taille $C+1$. \emph{Attention au sens de parcours} pour ne pas réutiliser deux fois
le même élément. Donner sa complexité en temps et en espace, et expliquer le terme
\emph{pseudo-polynomial}.

\Q Prouver la correction de \code{peut_partitionner} en énonçant un \emph{invariant} portant sur
\code{dp} après le traitement des $i$ premiers éléments, puis en le démontrant par récurrence sur
$i$. Justifier enfin la terminaison.

\clearpage
\section*{Problème II --- Arbres de Merkle et preuves d'inclusion}
\noindent\emph{Présentation.} Un \emph{arbre de Merkle} permet de résumer une grande collection de données par un unique
condensé, la \emph{racine}, et de prouver qu'une donnée en fait partie à l'aide d'un
\emph{certificat} de taille logarithmique. C'est un outil central des blockchains (preuves
d'inclusion, listes d'autorisation, stockage vérifiable). Dans tout le problème, les
\emph{feuilles} sont des entiers, et l'on combine deux n\oe uds par la fonction
$$\mathtt{parent}(a,b) = (31\,a + b) \bmod p, \qquad p = 1000003 \ (\text{premier}).$$

\noindent\textbf{Conventions.} Le langage est Python 3. Pour une liste \code{L}, \code{len(L)}
est sa longueur, \code{L[i]} son terme d'indice $i$ ($\ge 0$), \code{L[-1]} son dernier terme.
$\lceil x \rceil$ est la partie entière supérieure, $\log_2$ le logarithme binaire. Les questions
sont numérotées \textbf{en continu}.

\medskip
\subsection*{Partie I --- La fonction de combinaison}

\Q Écrire \code{parent(a, b)} conforme à la définition ci-dessus.

\Q Calculer à la main $\mathtt{parent}(1,2)$, $\mathtt{parent}(3,4)$ et
$\mathtt{parent}(33,97)$.

\Q On dit que \code{parent} est \emph{déterministe}. Expliquer pourquoi ce caractère déterministe
est indispensable pour que deux n\oe uds différents du réseau, partant des mêmes feuilles,
calculent la même racine.

\Q On admet, pour la suite, que \code{parent} est \emph{résistante aux collisions}~: on ne sait
pas exhiber deux couples distincts $(a,b)\neq(a',b')$ tels que $\mathtt{parent}(a,b)=
\mathtt{parent}(a',b')$. Rappeler pourquoi une fonction $h : \mathbb{N}^2 \to \{0,\dots,p-1\}$ ne
peut pas être \emph{injective}, et donc pourquoi la résistance aux collisions est une hypothèse
\emph{calculatoire} (« difficile à violer ») et non une impossibilité mathématique.

\subsection*{Partie II --- Construction de la racine}

La racine se calcule par niveaux~: tant qu'il reste plus d'un n\oe ud, si le nombre de n\oe uds
est impair on \textbf{duplique le dernier}, puis on remplace la liste par celle des
\code{parent} de couples consécutifs.

\Q Écrire \code{racine_merkle(feuilles)} (liste non vide) renvoyant la racine.

\Q Dérouler le calcul de \code{racine_merkle([1, 2, 3, 4])} niveau par niveau et donner la valeur
de la racine.

\Q Faire de même pour \code{racine_merkle([1, 2, 3])} en explicitant l'étape de duplication.

\Q Soit $n$ le nombre de feuilles. En notant $t_0=n$ et $t_{k+1}=\lceil t_k/2\rceil$ la taille du
niveau $k+1$, montrer que le nombre de niveaux est $O(\log n)$ et que le nombre total d'appels à
\code{parent} est un $O(n)$.

\subsection*{Partie III --- Preuves d'inclusion}

\Q Écrire \code{preuve_inclusion(feuilles, i)} renvoyant la liste des couples
$(\text{frère}, \text{côté})$, avec côté valant \code{'L'} ou \code{'R'}, rencontrés en remontant
de la feuille d'indice $i$ jusqu'à la racine.

\Q Écrire \code{verifier(feuille, preuve, racine)} qui recalcule la racine depuis la feuille et
la preuve et renvoie un booléen. Vérifier sur l'exemple $[1,2,3,4]$, $i=0$.

\Q Montrer que la taille d'une preuve d'inclusion est le nombre de niveaux, soit $O(\log n)$, et
commenter l'intérêt (comparer au coût de transmettre toutes les feuilles).

\Q \textbf{Soundness.} Démontrer par récurrence sur la hauteur $h$ de l'arbre que, sous
l'hypothèse de résistance aux collisions, si \code{verifier(feuille, preuve, racine)} renvoie
\code{True} alors \code{feuille} figure effectivement parmi les feuilles ayant produit
\code{racine}.

\subsection*{Partie IV --- Hauteur et mise à jour}

\Q Montrer que la hauteur de l'arbre (nombre de niveaux au-dessus des feuilles) vaut
$H(n) = \lceil \log_2 n \rceil$ pour $n \ge 1$ (avec $H(1)=0$). Donner $H(n)$ pour
$n \in \{1,2,3,4,5,8,9\}$.

\Q Lorsqu'une seule feuille change de valeur, quels n\oe uds de l'arbre doivent être recalculés~?
En déduire qu'une mise à jour de la racine coûte $O(\log n)$ appels à \code{parent} (et non
$O(n)$). Écrire \code{maj_racine(feuilles, i, v)} qui renvoie la nouvelle racine lorsque la feuille
d'indice $i$ prend la valeur \code{v}, en \textbf{réutilisant} \code{racine_merkle}. \emph{On ne
demande pas ici l'implémentation optimale en $O(\log n)$, mais on justifiera la borne.}

\clearpage
\section*{Problème III --- Le problème de partition et la programmation dynamique}
\noindent\emph{Présentation.} Le \emph{problème de partition} consiste, étant donné une liste $S$ d'entiers strictement
positifs, à décider si l'on peut la séparer en deux sous-ensembles de même somme. Il apparaît dans
les preuves à divulgation nulle de connaissance (chapitre « Zero-Knowledge Proofs » de
\emph{Mastering Ethereum}) comme exemple canonique de problème \emph{facile à vérifier} mais
\emph{difficile à résoudre}. On note $T = \sum_{x\in S} x$ la somme totale.

\noindent\textbf{Conventions.} Langage Python 3~; \code{sum(S)} est la somme des éléments,
\code{len(S)} leur nombre. Le module \code{itertools} fournit \code{combinations(S, r)} qui énumère
les sous-ensembles de $S$ à $r$ éléments. Questions numérotées \textbf{en continu}.

\medskip
\subsection*{Partie I --- Recherche exhaustive}

\Q Écrire \code{demi_somme(S)} qui renvoie \code{None} si $T=\mathtt{sum(S)}$ est impair, et
$T/2$ (entier) sinon.

\Q Écrire \code{partition_exhaustive(S)} qui décide, en énumérant tous les sous-ensembles à l'aide
de \code{itertools.combinations}, s'il existe un sous-ensemble de somme $T/2$. La fonction renvoie
un booléen. On renverra \code{False} si $T$ est impair.

\Q Combien de sous-ensembles possède un ensemble à $n$ éléments~? En déduire que
\code{partition_exhaustive} a une complexité temporelle en $\Theta(2^{n})$ dans le pire des cas,
et expliquer pourquoi cette approche est inutilisable pour $n$ grand.

\subsection*{Partie II --- Le problème de décision Subset-Sum}

On suppose désormais $T$ paire et on pose $C = T/2$. Le \emph{problème de décision} Subset-Sum
demande~: « existe-t-il un sous-ensemble de $S$ de somme exactement $C$~? »

\Q Décrire un \emph{certificat} (une preuve courte) permettant de convaincre en temps polynomial
qu'une instance est positive, et l'algorithme de \emph{vérification} associé. Donner sa complexité.

\Q En déduire que Subset-Sum appartient à la classe \textbf{NP}. \emph{On rappelle qu'un problème
de décision est dans NP s'il admet un certificat vérifiable en temps polynomial.} On mentionnera,
sans le démontrer, que Subset-Sum est en fait NP-complet.

\subsection*{Partie III --- Récurrence et table de programmation dynamique}

On définit, pour $0 \le i \le n$ (avec $n=\mathtt{len(S)}$) et $0 \le c \le C$, le booléen
$$A[i][c] = \text{« il existe un sous-ensemble des } i \text{ premiers éléments de } S
\text{ de somme } c \text{ »}.$$

\Q Justifier la relation de récurrence
$$A[i][c] = A[i-1][c] \;\lor\; \bigl(c \ge S[i-1] \ \land\ A[i-1][c - S[i-1]]\bigr)
\quad (i \ge 1),$$
et préciser l'initialisation $A[0][\,\cdot\,]$.

\Q En déduire une fonction \code{subset_sum_2d(S, C)} qui construit la table $A$ (de taille
$(n+1)\times(C+1)$) et renvoie \code{A[n][C]}. Donner sa complexité en temps et en espace.

\subsection*{Partie IV --- Version optimisée en espace}

\Q Écrire \code{peut_partitionner(S)} qui, après avoir écarté le cas $T$ impair, décide
l'existence d'une partition équilibrée en n'utilisant qu'un tableau \code{dp} \textbf{à une
dimension} de taille $C+1$. \emph{Indication~:} parcourir $c$ de $C$ jusqu'à $x$ en
\textbf{décroissant} pour chaque élément $x$ de $S$.

\Q Expliquer précisément pourquoi le parcours \textbf{décroissant} de $c$ est nécessaire~: que se
passerait-il avec un parcours croissant~?

\Q Donner la complexité en temps et en espace de \code{peut_partitionner}, et expliquer le terme
\emph{pseudo-polynomial} (comparer $C$ à la taille de codage de l'entrée).

\subsection*{Partie V --- Reconstruction et correction}

\Q À partir de la table $A$ de la Partie III, écrire \code{sous_ensemble(S, C)} qui renvoie
\emph{un} sous-ensemble de somme $C$ (sous forme de liste), ou \code{None} s'il n'en existe pas.
\emph{Indication~:} remonter la table de $i=n$ à $1$~; l'élément $S[i-1]$ est pris ssi
$A[i-1][C]$ est faux (le retirer impose de l'utiliser).

\Q Énoncer un \emph{invariant} portant sur \code{dp} après le traitement des $i$ premiers éléments
dans \code{peut_partitionner}, puis le démontrer par récurrence sur $i$. Conclure sur la
correction, et justifier la terminaison.

\clearpage
\section*{Problème IV --- Finance décentralisée : teneurs de marché automatisés}
\noindent\emph{Présentation.} Un \emph{teneur de marché automatisé} (\emph{Automated Market Maker}, AMM) remplace le carnet
d'ordres par une \emph{réserve de liquidité} contenant deux jetons, en quantités $x$ et $y$. Le
prix est fixé par une formule et non par l'appariement d'ordres. Le modèle le plus répandu
(Uniswap V2) impose que le \emph{produit} $x\,y$ reste constant~: $x\,y = k$. Ce problème étudie
ce mécanisme, ses bornes, puis un modèle de \emph{prêt sur-collatéralisé}.

\noindent\textbf{Conventions.} Langage Python 3. Les réserves $x,y$ sont des réels
strictement positifs. On note $\Delta x > 0$ la quantité de jeton $X$ apportée à la réserve et
$\Delta y > 0$ la quantité de jeton $Y$ reçue en échange (échange \emph{sans frais}, sauf mention
contraire). Questions numérotées \textbf{en continu}.

\medskip
\subsection*{Partie I --- Prix et produit constant}

\Q Écrire \code{prix(x, y)} qui renvoie le prix d'une unité du jeton $X$ exprimé en jeton $Y$,
c'est-à-dire $y/x$. Vérifier que pour une réserve de $10$ ETH et $100$ USDC, le prix d'un ETH est
$10$ USDC.

\Q Un échange amène la réserve de $(x, y)$ à $(x + \Delta x,\, y - \Delta y)$. En imposant la
conservation du produit $(x+\Delta x)(y-\Delta y) = x\,y$, montrer que
$$\Delta y = \frac{y\,\Delta x}{x + \Delta x}.$$

\Q Écrire \code{sortie_amm(x, y, dx)} qui renvoie $\Delta y$. Vérifier que
\code{sortie_amm(100, 100, 100)} vaut $50$.

\subsection*{Partie II --- Bornes et monotonie (anti-drainage)}

\Q Montrer que pour tout $\Delta x > 0$ on a $\Delta y < y$~: on ne peut jamais vider entièrement
la réserve du jeton $Y$.

\Q Montrer que la fonction $\Delta x \mapsto \Delta y(\Delta x) = \dfrac{y\,\Delta x}{x+\Delta x}$
est \textbf{strictement croissante} sur $]0, +\infty[$ et que $\Delta y \to y$ lorsque
$\Delta x \to +\infty$. Interpréter en termes de \emph{glissement} (\emph{slippage})~: plus
l'échange est gros relativement à la réserve, plus le prix effectif se dégrade.

\Q Le \emph{prix effectif} d'un échange est $\Delta x / \Delta y$. L'exprimer en fonction de
$x, y, \Delta x$ et montrer qu'il est $> x/y$ (le prix « spot » inverse), avec égalité
\emph{à la limite} $\Delta x \to 0$. Conclure sur l'impossibilité de drainer la réserve à prix
constant.

\subsection*{Partie III --- Frais cumulés et série géométrique}

On modélise une suite d'échanges dont les volumes décroissent géométriquement~: le $i$-ième
échange ($i \ge 0$) porte un volume $v_0\,\rho^{\,i}$ avec $0 < \rho < 1$, et prélève des frais
proportionnels de taux $f$.

\Q Montrer que le total des frais perçus $\displaystyle \Phi = \sum_{i=0}^{+\infty} f\,v_0\,
\rho^{\,i}$ est fini et vaut $\dfrac{f\,v_0}{1-\rho}$.

\Q Application~: pour $f = 0{,}003$, $v_0 = 1000$ et $\rho = 0{,}9$, calculer $\Phi$.

\Q Écrire \code{frais_cumules(f, v0, rho, n)} qui calcule la somme \emph{partielle}
$\sum_{i=0}^{n-1} f\,v_0\,\rho^{\,i}$, et donner une majoration de l'erreur
$\Phi - \mathtt{frais\_cumules(f, v0, rho, n)}$ par une expression géométrique en $\rho^{\,n}$.

\subsection*{Partie IV --- Prêt sur-collatéralisé et intérêts composés}

Un emprunteur dépose un \emph{collatéral} de valeur $G$ et emprunte un principal $P$. Le taux de
collatéralisation exigé est $\tau\%$~: l'emprunt n'est autorisé que si $G \ge P\,\tau/100$. La dette
croît ensuite par intérêts composés~: après $n$ périodes au taux $r$ par période, la dette vaut
$P(1+r)^n$. La position est \emph{liquidée} dès que la dette dépasse le collatéral $G$.

\Q Écrire \code{collateral_requis(P, tau)} qui renvoie le collatéral minimal $\lfloor P\,\tau/100
\rfloor$ (division entière). Vérifier \code{collateral_requis(100, 150) == 150}.

\Q Écrire \code{dette(P, r, n)} renvoyant $P(1+r)^n$. Montrer que la dette dépasse strictement le
collatéral $G$ pour la première fois à la période
$$n^{\star} = \left\lfloor \frac{\ln(G/P)}{\ln(1+r)} \right\rfloor + 1
\qquad (\text{si } G > P).$$

\Q Application~: pour $P = 100$, $r = 0{,}10$ et $G = 150$, calculer $n^{\star}$ et vérifier
l'encadrement $\mathtt{dette}(100, 0.1, n^{\star}-1) \le 150 < \mathtt{dette}(100, 0.1, n^{\star})$.
Justifier la \textbf{terminaison} d'une boucle \code{while} qui incrémenterait $n$ tant que
$\mathtt{dette}(P, r, n) \le G$.

\end{document}
